\section{Intermediate singular values}\label{sec:intermediate}

We convert~\eqref{eq:averaged-normal-small-ball} into a lower bound for an intermediate singular value. Fix $z\in\C$ and write
\[
  A_z=M_n-z\sqrt d\,I_n.
\]
Let $N=n-m$ with
\[
  m=\floor{n/d^3}-1,
\]
and expose the first $N$ rows $R_1,\ldots,R_N$ of $A_z$.

For each $i\le N$, let
\[
  F_i=\spanv\set{\ol{R_j}:j\le N,\ j\ne i}
\]
and define the actual leave-one-row-out distance
\[
  D_i=\dist(\ol{R_i},F_i)
  =\dist\bigl(R_i,\spanv\set{R_j:j\le N,\ j\ne i}\bigr).
\]
If $\dim F_i<N-1$, set $E_i=\Ext_{N-1}(F_i)$; otherwise set $E_i=F_i$. Put
\[
  \wt D_i=\dist(\ol{R_i},E_i)\le D_i,
  \qquad
  Y_i=P_{E_i^\perp}G_i,
\]
where $G_i$ is an independent standard complex Gaussian. Conditional on $R_i$ and $E_i$,
\[
  R_i\cdot Y_i
  =\inner{Y_i}{\ol{R_i}}
  =\inner{G_i}{P_{E_i^\perp}\ol{R_i}}
\]
is a centered complex Gaussian with variance $\wt D_i^2$. Hence there are absolute $c,c_0>0$ such that
\begin{equation}\label{eq:Gaussian-distance-implication}
  \wt D_i\le c\rho
  \quad\Longrightarrow\quad
  \Pp_{G_i}\set{|R_i\cdot Y_i|\le\rho\given R_i,E_i}\ge c_0,
\end{equation}
where
\[
  \rho=e^{-C_{\isv}(z)d^{1/96}}
\]
for a sufficiently large shift-dependent constant $C_{\isv}(z)$.

Conditional on $E_i,Y_i$, the omitted row $R_i=X_i-z\sqrt d\,e_i$ remains an independent row slice. The diagonal shift changes only the center of the concentration function, so~\eqref{eq:averaged-normal-small-ball} gives
\begin{equation}\label{eq:row-normal-small-ball}
  \Pp\set{|R_i\cdot Y_i|\le\rho}\le Cd^{-3}+C/n.
\end{equation}
Combining~\eqref{eq:Gaussian-distance-implication} and~\eqref{eq:row-normal-small-ball},
\begin{equation}\label{eq:leave-one-out-distance-prob}
  \Pp\set{D_i\le c\rho}
  \le\Pp\set{\wt D_i\le c\rho}
  \le Cd^{-3}+C/n.
\end{equation}

For each fixed index $i$, the exposed set $\set{1,\ldots,N}\setminus\set i$ is deterministic, so the fixed-$I$ estimate in 
\cref{sec:normals} applies. We do not intersect these $N$ structural events and do not take a union bound over the leave-one-out indices. Instead, the per-index bound is summed only after taking expectations.

Let
\[
  B=\#\set{i\le N:D_i<c\rho}.
\]
Linearity of expectation, not independence, gives
\[
  \E B\le C\frac n{d^3}+C.
\]
Set
\[
  t=\ceil{\frac n{d^2}}.
\]
Markov's inequality yields
\begin{equation}\label{eq:bad-row-Markov}
  \Pp\set{B>t}
  =O\left(\frac1d+\frac{d^2}{n}\right)
  =O(1/d)=o(1).
\end{equation}

On $B\le t$, delete all bad rows and, if necessary, additional arbitrary rows so that the resulting matrix $H$ has $k=N-t$ rows. For every surviving row $h_i$,
\[
  \dist\bigl(h_i,\spanv\set{h_j:j\ne i}\bigr)\ge c\rho,
\]
because deleting rows only shrinks the comparison span. By the negative second-moment identity~\eqref{eq:negative-second-moment},
\begin{equation}\label{eq:NSM-H}
  \sum_{j=1}^ks_j(H)^{-2}\le\frac{k}{c^2\rho^2}.
\end{equation}
Let $j_0=N-2t=k-t$. At least $t+1$ singular values of $H$ are no larger than $s_{j_0}(H)$, and therefore
\[
  \frac{t+1}{s_{j_0}(H)^2}
  \le\sum_{j=1}^ks_j(H)^{-2}.
\]
Using~\eqref{eq:NSM-H},
\begin{equation}\label{eq:intermediate-H}
  s_{j_0}(H)
  \ge c\rho\sqrt{\frac{t+1}{k}}
  \ge\frac{c\rho}{d}.
\end{equation}
Adding the deleted rows replaces the common $n\times n$ column Gram matrix $H^*H$ by
\[
  H^*H+\sum_{u\in\mathcal D}R_u^*R_u,
\]
a positive-semidefinite increment. Hence the relevant singular values can only increase, and
\[
  s_{j_0}(A_z)\ge s_{j_0}(H).
\]
Define
\[
  q_0=m+2t.
\]
Then $j_0=n-q_0$ and
\begin{equation}\label{eq:q0-scale}
  q_0=O\left(\frac n{d^2}\right)
  =o\left(\frac n{d\log d}\right).
\end{equation}
Since the polynomial factor $d^{-1}$ is absorbed into the exponential,~\eqref{eq:intermediate-H} proves the desired estimate.

\begin{theorem}[Lower-band intermediate singular value]\label{thm:intermediate}
Fix $\sigma\in(0,1)$ and let $d$ satisfy~\eqref{eq:lower-band}. Let $M_n$ be the independent row-slice matrix. For every fixed $z\in\C$, there are constants $C_{\isv}(z),C'_{\isv}(z)>0$ and integers $q_0=q_0(n,z)$ with
\[
  q_0\le C'_{\isv}(z)\frac n{d^2}
\]
such that
\[
  \Pp\set{s_{n-q_0}(M_n-z\sqrt d\,I_n)
  \ge e^{-C_{\isv}(z)d^{1/96}}}=1-o(1).
\]
More precisely, the failure probability is $O(1/d)$ uniformly throughout~\eqref{eq:lower-band}; the term $d^2/n$ and the exponentially small net and witness errors are absorbed into this bound.
\end{theorem}

\begin{remark}
The choice
\[
  m=\floor{n/d^3}-1=\frac n{d^3}+O(1)
\]
is tuned to the LLTTY exponent $\gamma=1/288$:
\[
  (n/m)^\gamma=d^{1/96+o(1)},
  \qquad
  (m/n)^{\gamma/2}=d^{-1/192+o(1)}.
\]
The deletion scale $t=n/d^2$ is larger than the expected number $n/d^3$ of bad rows, but remains small enough for the lower-band logarithmic-tail budget.
\end{remark}
